\documentclass[10pt,a4paper]{amsart}
\usepackage[margin=19mm]{geometry}
\usepackage{amsmath,amssymb,mathtools}
\usepackage[scr=rsfs,cal=euler]{mathalfa}
\usepackage{xfrac}
\usepackage[unicode,hidelinks,bookmarks=false]{hyperref}
\newcommand{\ie}{i.e.\ }
\newcommand{\cf}{cf.\ }
\newcommand{\resp}{respectively\ }
\newcommand{\Subsection}{\subsection}
\providecommand{\nobreakdash}{\nobreak\hbox{-}}
\setlength{\emergencystretch}{2em}
\multlinegap0pt
\usepackage{bm}
\usepackage{bbm}
\usepackage{dsfont}
\usepackage{xspace}
\usepackage{cancel}
\usepackage{mathtools}
\usepackage{amsthm}
\makeatletter
\@ifundefined{theorem}{\newtheorem{theorem}{Theorem}}{}
\@ifundefined{proposition}{\newtheorem{proposition}[theorem]{Proposition}}{}
\makeatother
\DeclareMathOperator{\seg}{seg}
\newcommand{\ba}{\mathbin{\backslash}}
\newcommand{\calM}{\mathcal{M}}
\newcommand{\cog}[1]{\textbf{\textit #1}}
\newcommand{\ext}{\mathbin{\dagger}}
\newcommand{\gec}[1]{\mathcal{#1}}
\title[Polynomial invariants of cyclically ordered graphs]{Polynomial invariants of cyclically ordered graphs}
\author{Paul Bratch, M. N. Ellingham, Joanna A. Ellis-Monaghan, Iain Moffatt, Wout Moltmaker}
\date{Source: arXiv:2511.13302v1 (CC BY 4.0)}
\begin{document}
\maketitle

\noindent\emph{Source and licence.} Polynomial invariants of cyclically ordered graphs. Version arXiv:2511.13302v1 (CC BY 4.0), distributed under the Creative Commons Attribution 4.0 International licence, as recorded in the arXiv licence metadata for this exact version and shown on the abstract page https://arxiv.org/abs/2511.13302v1. Full rights notice: licenses/arxiv-2511.13302-CC-BY-4.0.txt.\par\smallskip

\noindent\textit{Editorial scope.} Counted unit: the Theorem whose source label is \texttt{thm:scg} in the article (source0.tex lines 746--749, source label \texttt{thm:scg}), delivered together with the article's own complete original proof of it (source0.tex lines 750--782, one \texttt{proof} environment of 33 lines). No mathematics is written, repaired or re-derived: statement and proof are the article's own text line for line. Required context retained uncounted: prop:2dag (lines 625--629); thm:mwd (lines 638--646); thm:rda (lines 675--681). Every definition, display and label of those lines is retained unchanged. Omitted from this package and not redistributed: 1--624, 630--637, 647--674, 682--745, 783--1805. Nothing inside a retained excerpt is omitted or altered: every retained body is delivered exactly as the article's lines are written, so no omission receipt of any kind accompanies this package. Citations: the retained excerpts carry no citation command at all, so no bibliography entry of the article is transcribed and no reference is redistributed. Reference adaptation: none was needed and none was made; every reference inside the delivered unit resolves inside it, so no unresolved reference or citation is produced. Printed slips kept literal: no printed word, symbol, hypothesis or number of the counted statement or of its proof was altered. Layout and class: the article's own class is \texttt{amsart} with packages including mathtools, amssymb, amsmath, graphicx, pinlabel, pdfpages, lineno, xspace, comment, accents, subcaption, hyperref, xcolor, fancyhdr; the wrapper is the lane's pinned \texttt{amsart} editorial wrapper at 10pt on A4, which restates the theorem-like environments the retained text needs and adds the article's own macro definitions that the retained text needs (\texttt{\textbackslash ba}, \texttt{\textbackslash calM}, \texttt{\textbackslash cog}, \texttt{\textbackslash ext}, \texttt{\textbackslash gec}, \texttt{\textbackslash seg}), with extra wrapper packages (bm, bbm, dsfont, xspace, cancel, mathtools). The delivered numbering is the unit's own. Assets: no retained span contains a figure environment, a graphics-inclusion command, a PDF-page inclusion, a table or a TikZ picture; no image file is copied into this package, the package declares no asset dependency and no image file is redistributed with it. Licence: version arXiv:2511.13302v1 of this article is distributed under the Creative Commons Attribution 4.0 International licence, as recorded in the arXiv licence metadata for this exact version and shown on the abstract page https://arxiv.org/abs/2511.13302v1. A full rights notice is delivered at licenses/arxiv-2511.13302-CC-BY-4.0.txt. Characters: every retained excerpt is pure ASCII, so the delivered engine, XeLaTeX with its default Latin Modern text font, reports no missing character.
\begin{proposition}\label{prop:2dag}
Let $\gec{G}=(V,E)$ be a generalised gec,  $\mathcal{E}$ be its set of e-edges, and $e,f\in \mathcal{E}$. Then 
$(\gec{G}\ext e)\ext f = (\gec{G}\ext f)\ext e$ and 
$(\gec{G}\ext e)\ba f = (\gec{G}\ba f)\ext e$. 
\end{proposition}

\begin{theorem}\label{thm:mwd}
Let $\gec{G}=(V,E)$ be a generalised gec and $\mathcal{E}$ be its set of e-edges. Then
\[
\calM(\gec{G};x,y) 
= \sum_{\mathcal{A} \subseteq \mathcal{E}} y^{|\mathcal{A}|}x^{k(\gec{G}\ext \mathcal{A} \setminus (\mathcal{E}-\mathcal{A}))}
= \sum_{\mathcal{A} \subseteq \mathcal{E}} y^{|\mathcal{A}|}x^{k_{\textup{v}}(\gec{G}\ext \mathcal{A} )},
\]
where $k_{\textup{v}}(\gec{G})$ means the number of v-components of a generalised gec $\gec{G}$.
\end{theorem}

\begin{theorem}\label{thm:rda}
Let $\cog{G}=(V,E)$ be a cog and $\gec{G}$ be its corresponding gec. Then
\[
 \calM(\gec{G};x, y) = \sum_{B\subseteq E}   y^{|E|-|B|} x^{\seg(B)+\iota(\cog{G})}, 
 \]
where $\seg(B)$ is the number of $B$-segments in the cog $\cog{G}$, and $\iota(\cog{G})$ the number of isolated vertices in $\cog{G}$.
\end{theorem}

\begin{theorem}\label{thm:scg}
Let $\cog{G}=(V,E)$ be a cog and $\gec{G}$ be its corresponding gec. In addition let $D$ and $X$ be disjoint subsets of $E$, and let $\gec{D}$ and $\gec{X}$ be the corresponding sets of $e$-edges of $\gec{G}$.  Then
\[ M(\cog{G},D,X;x, y) =  \calM(\gec{G}\ba \gec{D}\ext \gec{X}  ; x,y).\]
\end{theorem}

\begin{proof}
The proof follows that of Theorem~\ref{thm:rda}.
Let $G$, $\gec{G}$,  $D$, $X$, $\gec{D}$ and $\gec{X}$ be as in the theorem statement, and let $\mathcal{E}$ be the set of e-edges in $\gec{G}$, so $\mathcal{E}-(\gec{D}\cup \gec{X})$ is the set of e-edges in $\gec{G}\ba \gec{D}\ext \gec{X}$. 
Reindexing the sum in Theorem~\ref{thm:mwd} and using Proposition~\ref{prop:2dag} gives  
\begin{align*}
\calM(\gec{G}\ba \gec{D}\ext \gec{X};x,y) 
&= \sum_{\mathcal{B} \subseteq \mathcal{E}-(\gec{D}\cup \gec{X})} y^{|\mathcal{E}-(\gec{D}\cup \gec{X})|-|\mathcal{B}|}x^{k((\gec{G}\ba \gec{D}\ext \gec{X})\ba \mathcal{B} \ext (\mathcal{E}-(\gec{D}\cup \gec{X})-\mathcal{B}))}
\\
&=
\sum_{\mathcal{B} \subseteq \mathcal{E}-(\gec{D}\cup \gec{X})} 
y^{|\mathcal{E}-(\mathcal{B}\cup \gec{D}\cup \gec{X})|}
x^{k(\gec{G}\ba (\gec{D}\cup \gec{B}) \ext (\gec{E}-(\gec{D}\cup \gec{B})))}
.
\end{align*}

Fix some $\mathcal{B} \subseteq \mathcal{E}-(\gec{D}\cup \gec{X})$ and let $B\subseteq E$ be the corresponding set of edges in the cog $\cog{G}$. 
The gec
$\gec{G}\ba (\gec{D}\cup\mathcal{B}) \ext (\mathcal{E}-(\gec{D}\cup\mathcal{B}))$ can be obtained by taking the union of all of the v-cycles in $\gec{G}$ then deleting each vertex that was not an end of an edge in $\gec{D}\cup\mathcal{B}$. 
Thus, the contribution of a cycle of v-edges $C_v$ to $k(\gec{G}\ba (\gec{D}\cup\mathcal{B}) \ext (\mathcal{E}-(\gec{D}\cup\mathcal{B})))$ can be obtained by deleting all the vertices in $C_v$ that were not an end of an edge in $\gec{D}\cup\mathcal{B}$. When $C_v$ is not a free loop, this contribution is easily seen to equal $\seg(v,D\cup B)$ where $v$ is the vertex in the cog $\cog{G}$ corresponding to $C_v$. 
When $C_v$ is a free loop it contributes one to the number of connected components, and the free loop corresponds to a unique isolated vertex in
$\cog{G}$.
Thus, 
$k(\gec{G}\ba (\gec{D}\cup\mathcal{B}) \ext (\mathcal{E}-(\gec{D}\cup\mathcal{B})))$ 
is equal to 
$\seg(B\cup D)+\iota(\cog{G})$. 
Additionally,  $|\mathcal{E}-(\mathcal{B}\cup \gec{D}\cup \gec{X})| = |E|-|B\cup D\cup X|$.
Thus, for corresponding sets $\gec{B}$ and $B$ we have  
$y^{|\mathcal{E}-(\mathcal{B}\cup \gec{D}\cup \gec{X})|}
x^{k(\gec{G}\ba (\gec{D}\cup \gec{B}) \ext (\gec{E}-(\gec{D}\cup \gec{B})))}
=
 y^{|E|-|B\cup D\cup X|} x^{\seg(B\cup D)+\iota(\cog{G})}$ and
the result follows.
\end{proof}

\end{document}
